\section{Keep the six completed interfaces explicit}
\texttt{qoi\_bound} accepts packed dimensions in \texttt{a0}, width in the low
thirty-two bits and height in the high thirty-two. Valid dimensions return zero in
\texttt{v0} and \(22+5wh\) in \texttt{v1}; invalid dimensions return one and zero.
This helper needs no caller pointer.

\texttt{qoi\_validate} accepts file bytes, unsigned count and metadata pointer in
\texttt{a0,a1,a2}. A valid stream publishes sixteen little-endian metadata bytes:
u32 width at zero, u32 height at four, u32 channels at eight and u32 colorspace at twelve.
The metadata pointer is eight-byte aligned. Success returns pixel count in \texttt{v1};
failure changes no metadata byte.

\texttt{qoi\_decode} accepts file bytes, unsigned count, output word pointer and unsigned
pixel capacity in \texttt{a0} through \texttt{a3}. Its output is four-byte aligned.
Success returns packed dimensions. It validates the stream before capacity and writes
no output on failure. \texttt{qoi\_encode} receives input RGBA words in \texttt{a0},
a pointer to the same sixteen-byte descriptor in \texttt{a1}, destination bytes in
\texttt{a2} and unsigned byte capacity in \texttt{a3}. It retains the supplied channels
and colorspace fields in the header and returns the actual encoded byte length.
This metadata preservation matters when an input file describes linear components.

For these memory routines, zero status means success, one means malformed data or
invalid descriptor fields, and two means a count or capacity bound. Failures return
zero count. Validator and decoder reject counts above 20502 before pointer access;
counts below twenty-three are malformed before pointer access. Encoder descriptor
validation precedes capacity, which precedes pixel reads and output writes.

\texttt{qoi\_load} receives a path, byte buffer, unsigned capacity and metadata pointer.
Require capacity from one through 20502 before pointer or file access. Read through EOF,
close, then validate and publish metadata. Return the loaded byte count. The byte buffer
may retain a prefix on failure; metadata must remain unchanged. \texttt{qoi\_save}
receives a path, bytes and unsigned count. Validate before opening in truncate mode,
then handle short writes and close on every return. Success returns the saved count.
Its count guards match the validator: above 20502 returns two and below twenty-three
returns one, before pointer or path access.
An open/read/write failure returns three and zero. A failed write may leave a partial file.
The current syscall's extra readable guard byte remains outside the selected file count.

Retain full-width pointers, counts, capacities, saved registers and return addresses.
Preserve \texttt{s0} through \texttt{s7}, \texttt{fp} and \texttt{sp}. Produce no console
I/O. Byte pointers may be unaligned; output objects are disjoint from their inputs.
Selected input intervals may be read, and only the declared result interval may be
written among caller objects. Private cache and frame storage belong to the invocation.
